\section{Анализ предметной области}
\subsection{Регулярные выражения}

Регулярные выражения (\textit{regex} -- от англ. \textit{Regular Expressions}) -- это формальный язык, используемый для поиска и осуществления манипуляций с фрагментами текста, основанный на использовании метасимволов (англ. \textit{wildcards}). Для поиска используется шаблон регулярного выражения -- строка-образец, состоящая из символов и метасимволов и позволяющая анализировать входящий текст удобным образом. Для манипуляций с текстом дополнительно задаётся строка замены, которая также может содержать в себе специальные символы.

Регулярные выражения широко используются в различных областях программирования и обработки данных, такими как:
\begin{itemize}
	\item текстовые редакторы и IDE (расширенный поиск и подстановка текста на основе шаблонов);
	\item валидация данных (проверка формата вводимых данных, таких как адреса электронной почты, номера телефонов, почтовые индексы, гиперссылки);
	\item парсинг данных (извлечение информации из неструктурированных текстов, таких как HTML, XML или лог-файлы);
	\item обработка данных в базах данных;
	\item фильтрация спама в системах почты.
\end{itemize}

Истоки регулярных выражений лежат в теории автоматов и формальных языков. Эти области изучают вычислительные модели (автоматы) и способы описания и классификации формальных языков. В 1940-x гг. Уоррен Маккалак и Уолтер Питтс описали нервную систему, используя простой автомат в качестве модели нейрона. Математик Стивен Клини позже описал эти модели, используя свою систему математических обозначений, названную «регулярные множества». Кен Томпсон встроил их в редактор QED, а затем в редактор ed под UNIX. С этого времени регулярные выражения стали широко использоваться в Unix и Unix-подобных утилитах, например: expr, awk, Emacs, vi, lex и Perl.

Регулярные выражения в Perl и Tcl происходят от реализации, написанной Генри Спенсером. Филип Хазэль разработал библиотеку pcre (англ. Perl Compatible Regular Expressions -- Perl-совместимые регулярные выражения), которая используется во многих современных инструментах, таких как PHP и Apache. 

Существует несколько разновидностей регулярных выражений:
\begin{enumerate}
	\item Регулярные выражения POSIX -- стандарт, используемый в Unix и Unix-подобных системах. Делится на два типа: <<базовые>> (BRE) и <<расширенные>> (ERE).
	\item Регулярные выражения, совместимые с Perl (PCRE) -- расширенная версия регулярных выражений, используемая в Perl и поддерживаемая многими языками, такими как PHP, R и JavaScript.
	\item Регулярные выражения JavaScript -- подобны PCRE, поддерживают флаги, такие как g (глобальный поиск) и i (игнорирование регистра).
	\item Шаблоны Lua -- упрощённая версия регулярных выражений, используемая в языке Lua. Имеет менее мощный синтаксис, чем в других языках, не поддерживает некоторые сложные конструкции.
\end{enumerate}

\subsection{Методы реализации обработчика регулярных выражений}

В теории и практике компиляторов и обработки формальных языков сложились три основных подхода к реализации механизма сопоставления регулярных выражений с текстом: интерпретация синтаксического дерева (дерева разбора), симуляция конечного автомата и рекурсивный спуск с возвратом.

\subsubsection{Интерпретация синтаксического дерева}

При данном подходе шаблон регулярного выражения сначала разбирается синтаксическим анализатором (парсером) и представляется в виде абстрактного синтаксического дерева (АСД, от англ. \textit{Abstract Syntax Tree}, AST) -- иерархической структуры данных, в которой каждый внутренний узел соответствует операции (конкатенация, альтернация, квантификатор), а листовые узлы -- атомарным элементам шаблона (конкретным символам или символьным классам). Сопоставление с текстом выполняется путём рекурсивного обхода этого дерева: для каждого узла определяется, может ли соответствующий фрагмент текста быть сопоставлен с данной ветвью дерева.

Достоинствами данного подхода являются концептуальная простота реализации и близость структуры данных к грамматике языка регулярных выражений, что упрощает поддержку расширенных возможностей: именованных групп, условных конструкций, атомарных групп. Помимо этого, данный подход удобен для дальнейшей трансляции -- дерево разбора служит естественным промежуточным представлением для последующей генерации кода автомата или байт-кода виртуальной машины.

Среди недостатков следует отметить, что наивная интерпретация дерева с возвратами (без мемоизации) обладает экспоненциальной сложностью в худшем случае и может быть подвержена атакам типа ReDoS (\textit{Regular Expression Denial of Service}). Эффективность работы напрямую зависит от применяемых оптимизаций дерева и стратегии обхода.

\subsubsection{Симуляция конечного автомата}

Данный подход основан на классическом результате теории формальных языков: каждому регулярному выражению соответствует конечный автомат, допускающий тот же язык. Реализация включает два этапа. На первом этапе шаблон регулярного выражения транслируется в недетерминированный конечный автомат (НКА) с использованием алгоритма Томпсона, который строит НКА в виде ориентированного графа переходов за время и пространство, линейное по размеру шаблона. На втором этапе выполняется симуляция НКА на входной строке: вместо перебора с возвратами одновременно отслеживается множество активных состояний, в которых может находиться автомат после обработки каждого символа.

Альтернативой является предварительная детерминизация: НКА преобразуется в детерминированный конечный автомат (ДКА) посредством алгоритма построения подмножеств, после чего сопоставление с текстом выполняется за строго линейное время относительно длины входной строки. Детерминизация может производиться заранее (статически) или лениво -- в процессе сопоставления (алгоритм Хопкрофта--Карпа, используемый, в частности, в библиотеке RE2 и пакете \texttt{regexp} языка Go).

Главным достоинством автоматного подхода является гарантированная линейная сложность сопоставления относительно длины входной строки для всего класса регулярных языков, что исключает возможность ReDoS. Это делает подход наиболее предпочтительным для применений, критичных к безопасности.

К ограничениям следует отнести сложность поддержки ряда расширений, выходящих за рамки регулярных языков: обратных ссылок (\textit{backreferences}), предпросмотров (\textit{lookahead/lookbehind}) и условных конструкций. Кроме того, полная детерминизация может приводить к экспоненциальному росту числа состояний ДКА (экспоненциальный взрыв), что делает хранение полного ДКА нецелесообразным для сложных шаблонов; именно поэтому в практических реализациях применяется ленивая детерминизация с ограничением кэша.

\subsubsection{Рекурсивный спуск с возвратами}

Рекурсивный спуск с возвратами (\textit{backtracking}) -- наиболее распространённый подход в практических реализациях (Perl, PCRE, Python \texttt{re}, большинство языков программирования). В этом подходе сопоставление выполняется прямым рекурсивным обходом шаблона параллельно с обходом входной строки: при неудачном сопоставлении в некоторой точке выполняется возврат к последней точке выбора и испытывается альтернативная ветвь.

Несмотря на экспоненциальную сложность в худшем случае, данный подход на практике эффективен для большинства реальных шаблонов. Его ключевым преимуществом является возможность естественной реализации расширений, нарушающих регулярность: обратных ссылок, атомарных групп, посессивных (сверхжадных) квантификаторов, предпросмотров и ретроспективных проверок, именованных групп захвата. Перечисленные конструкции недоступны или труднодостижимы в чистом автоматном подходе.

К недостаткам относятся: уязвимость к ReDoS при специально подобранных шаблонах и входных строках, а также зависимость производительности от порядка альтернатив в шаблоне. Ряд реализаций (PCRE2, Oniguruma) частично устраняют эти проблемы с помощью оптимизаций: мемоизации промежуточных результатов, посессивных квантификаторов и атомарных групп.

\subsubsection{Сравнительный анализ подходов}

Перечисленные подходы не являются взаимоисключающими: многие промышленные библиотеки сочетают их. Так, можно построить АСД на этапе разбора шаблона, затем транслировать его в НКА для основного сопоставления, а для конструкций вне класса регулярных языков переключаться на механизм с возвратами. Выбор подхода определяется требованиями к поддерживаемому синтаксису, гарантиям производительности и сложности реализации.

Разница в производительности между подходами на патологических шаблонах может быть колоссальной. Классический пример -- шаблон \verb`(a+)+b`, применённый к строке из $n$ символов \verb`a` без завершающего \verb`b`: алгоритм с возвратами совершает порядка $2^n$ проверок (например, при
$n = 30$ это свыше миллиарда операций и время выполнения в несколько минут), тогда как симуляция НКА обрабатывает ту же строку за $O(n \cdot m)$ шагов и завершается мгновенно. На реальных данных разница обычно менее выражена, однако на специально подобранных входных данных может достигать от тысячи до миллиона раз и более -- именно это используется в атаках ReDoS.

В таблице~\ref{tab:regex_approaches} приведено сравнение рассмотренных подходов по ключевым критериям.

\begin{xltabular}{\textwidth}{|C{3.8cm}|T|T|T|}
	\caption{Сравнение методов реализации обработчика регулярных выражений}
	\label{tab:regex_approaches} \\
	\hline
	\makecell{Критерий} & \makecell{АСД\\(интерпретация\\дерева)} & \makecell{Симуляция КА\\(НКА/ДКА)} & \makecell{Возврат\\(backtracking)} \\
	\hline
	\finishhead
	Сложность сопоставления & $O(2^n)$ в худшем случае & $O(n \cdot m)$ или $O(m)$ & $O(2^n)$ в худшем случае \\
	\hline
	Потребляемая память & $O(n)$ для стека рекурсии & $O(2^n)$ для ДКА в худшем случае; $O(n \cdot m)$ для НКА & $O(n)$ для стека вызовов \\
	\hline
	Обратные ссылки & Поддерживаются & Возможны на НКА, не поддерживаются на ДКА & Поддерживаются \\
	\hline
	Предпросмотры & Поддерживаются & Затруднены & Поддерживаются \\
	\hline
	Устойчивость к ReDoS & Низкая & Высокая (линейная) & Низкая \\
	\hline
	Сложность реализации & Средняя & Высокая & Ниже среднего \\
	\hline
	Типичные реализации & Промежуточное представление в PCRE, RE2, Rust \texttt{regex} & RE2, Go \texttt{regexp} & PCRE, Python, Perl \\
	\hline
\end{xltabular}

\subsection{Язык Си и C API}

C -- процедурный язык программирования общего назначения, в котором программы описываются наборами функций, оперирующих над данными через явные типы, указатели и структуры. Язык спроектирован так, чтобы давать программисту прямой и предсказуемый контроль над памятью и аппаратурой при минимальной абстракции над машиной.

C был создан в начале 1970-х Деннисом Ритчи в Bell Labs как развитие языка B и инструмент для разработки операционной системы UNIX. В ранних версиях язык уже поддерживал понятия типов, структур и побайтной работы с памятью, что сделало его удобным для системного программирования. Книга Кернигана и Ритчи <<The C Programming Language>> (1978) закрепила синтаксис и стиль, став практическим эталоном языка. В 1989 году появился стандарт ANSI C (C89/C90), затем шли ревизии C99, C11, C17/C18 и значительное обновление C23 -- все они добавляли новые возможности, уточняли поведение и расширяли стандартную библиотеку, сохраняя при этом совместимость в широкой степени.

Дизайн C опирается на несколько ключевых идей: минимализм синтаксиса, предсказуемость выполнения и высокая эффективность. Язык предоставляет минимальный набор строительных блоков -- переменные, выражения, управляющие конструкции, функции, указатели и структуры -- и поручает программисту управление ресурсами (выделение и освобождение памяти, контроль границ буферов). Такая философия делает C особенно пригодным там, где важны скорость, компактность и контроль над аппаратной частью.

C широко применяют в областях, где требуется низкоуровневое управление и высокая производительность. Это реализация ядра операционных систем и системных утилит, написание драйверов устройств и сетевых стеков, разработка встроенного ПО для микроконтроллеров и реального времени, создание высокопроизводительных серверных компонентов и баз данных, реализация частей компиляторов и рантаймов для других языков (например, ядро CPython), а также оптимизированные библиотеки для численных вычислений, графики и игровых движков. Благодаря стабильности ABI и простоте компоновки, многие библиотеки предоставляют C-интерфейс даже если их внутренняя реализация написана на других языках.

C API -- это распространенный способ сделать библиотеку доступной внешним программам: набор заголовочных файлов с декларациями, соответствующая бинарная библиотека и ясные правила владения ресурсами. В C API обычно экспортируют функции с единым префиксом (чтобы избежать конфликтов имен), используют инкапсуляцию с помощью непрозрачных указателей (opaque pointers) для скрытия внутренней реализации, возвращают коды ошибок и документируют, кто отвечает за выделение и освобождение памяти. Другие важные аспекты -- четкое версионирование ABI, требования по потокобезопасности и совместимость с C++ через \verb`extern "C"`.

Типовая схема небольшого C API выглядит как набор функций для создания и уничтожения контекста, операций над этим контекстом и вспомогательных функций для управления выделенной библиотекой памятью. Контракты функций должны быть явно задокументированы: допустимые входные значения, поведение при ошибках, правила владения результирующими ресурсами.

При проектировании C API важно установить однозначные правила владения памятью, избегать утечек и неопределенного поведения, инкапсулировать реализацию за опаковыми структурами и сохранять ABI-совместимость при развитии библиотеки. Для практической безопасности следует предусмотреть проверки размеров буферов, возвращать информативные коды ошибок и документировать, какие функции потокобезопасны.


\subsection{Фреймворк Qt}

Qt -- кроссплатформенный фреймворк для разработки приложений с графическим интерфейсом и сопутствующих подсистем (сеть, файловые операции, сериализация, многопоточность и т.д.), в основе которого лежит набор библиотек на C++. Qt предоставляет высокоуровневые абстракции для создания нативно выглядящих приложений на настольных системах, мобильных платформах и встраиваемых устройствах.

Это одна из наиболее зрелых и широко используемых кроссплатформенных библиотек для GUI-разработки: проект начат в начале 1990-х Хаавардом Нордом и Эйриком Чамбе-Энгом, впоследствии основавших компанию Trolltech, с тех пор прошел путь от коммерческой реализации к модели с открытым исходным кодом и коммерческой поддержкой. Qt стандартизировал подход к декларативному описанию интерфейсов (Qt Quick, QML) и расширил традиционные C++-возможности собственными механизмами метаобъектов (signals/slots, свойства, динамическая информация о типах).

Qt строится вокруг идеи модульности: рамки фреймворка разбиты на независимые модули -- Qt Core, Qt GUI, Qt Widgets, Qt Quick, Qt Network, Qt SQL и другие -- каждый из которых решает свой набор задач и предоставляет свой набор типов и сервисов. 

В центре архитектуры лежит объектная модель на C++ с расширенной метаобъектной системой. Она реализует механизм сигналов и слотов для гибкой коммуникации между объектами и поддерживает свойства, которые легко интегрируются с QML и системой метаданных. 

Qt предоставляет единый, кроссплатформенный API, благодаря чему приложения ведут себя и выглядят нативно на различных операционных системах -- Windows, macOS, Linux, iOS и Android -- и при этом могут использовать системные сервисы каждой платформы через единый интерфейс.
 
Для описания современных, декларативных интерфейсов Qt предлагает QML и движок Qt Quick: QML -- это язык разметки с поддержкой анимаций и встроенного JavaScript, который позволяет быстро проектировать UI и при необходимости связывать его с производительным C++‑бэкендом.

Qt также сопровождается развитой экосистемой инструментов: интегрированной средой разработки Qt Creator и визуальным редактором Qt Designer, официальными средствами сборки (qmake и поддержкой CMake), обширной документацией и множеством примеров, что облегчает разработку и распространение приложений на базе фреймворка.

Qt активно применяется при разработке графического интерфейса программ. Области применения фреймворка включают в себя:

\begin{itemize}
	\item десктопные приложения с богатым GUI (офисные программы, IDE, визуальные редакторы);
	\item встраиваемые и промышленные интерфейсы управления (панели приборов, терминалы);
	\item мобильные приложения и гибридные решения (Qt поддерживает сборку под iOS/Android);
	\item приложения с требованиями к визуализации и мультимедиа (CAD, графические редакторы, специализированные визуализаторы);
	\item инструменты разработки и среды исполнения (редакторы, отладчики, средства моделирования);
	\item сетевые и серверные компоненты (используя Qt Network и асинхронные механизмы).
\end{itemize}


 
